\documentclass[14pt]{beamer}
\usepackage{pepslides}
\title{Ray theory}
\lecturenumber{17}
\hypersetup{pdftitle={Lecture 17: Ray theory}}

\begin{document}

\titleslide

\begin{frame}[t]
\heading{From last time\ldots}
Plane wave ansatz:
\begin{eqn}
u_{i}(\mathbf{x},t)=f(t-p_{j}x_{j})a_{i}.
\end{eqn}
Christoffel equation for plane waves:
\begin{eqn}
\begin{gathered}
{}[\bm{\Gamma}(\mathbf{p})-\mathbf{I}]\mathbf{a}=\mathbf{0},\\[\eqgap]
\Gamma_{ik}(\mathbf{p})=\frac{1}{\rho}A_{ijkl}p_{j}p_{l}.
\end{gathered}
\end{eqn}
Definition of the slowness surface:
\begin{eqn}
\det[\bm{\Gamma}(\mathbf{p})-\mathbf{I}]=0.
\end{eqn}
\end{frame}

\begin{frame}
\centering
\includegraphics[height=0.55\paperheight]{newton_prism_engraving.jpg}
\begin{eqn}
\begin{gathered}
\nabla\times\mathbf{E}+\frac{\partial\mathbf{B}}{\partial t}=\mathbf{0},\quad\nabla\cdot\mathbf{B}=0,\\[\eqgap]
\nabla\times\mathbf{H}-\frac{\partial\mathbf{D}}{\partial t}=\mathbf{J},\quad\nabla\cdot\mathbf{D}=\rho.
\end{gathered}
\end{eqn}
\end{frame}

\begin{frame}
\begin{columns}[b]
\column{0.48\textwidth}
\centering
Quantum mechanics
\column{0.48\textwidth}
\centering
Classical mechanics
\end{columns}
\begin{columns}[c]
\column{0.48\textwidth}
\begin{eqn}
\ii\hbar\frac{\partial\Psi}{\partial t}=H\Psi
\end{eqn}
\column{0.48\textwidth}
\begin{eqn}
\begin{gathered}
\frac{\ddns\mathbf{q}}{\ddns t}=\frac{\partial H}{\partial\mathbf{p}},\\[\eqgap]
\frac{\ddns\mathbf{p}}{\ddns t}=-\frac{\partial H}{\partial\mathbf{q}}
\end{gathered}
\end{eqn}
\end{columns}
\end{frame}

\sectionslide{Length scales of heterogeneity}

\begin{frame}
\centering
\animation[height=0.6\paperheight]{L5F1.pdf}{L5F1.webm}

\raggedright
Scattering regime, with $\lambda$ the dominant wavelength and $L$ the length scale of the heterogeneity:
\begin{eqn}
\frac{L}{\lambda}\approx1.
\end{eqn}
\end{frame}

\begin{frame}
\centering
\animation[height=0.6\paperheight]{L5F2.pdf}{L5F2.webm}

\raggedright
Geometric optics regime:
\begin{eqn}
\frac{L}{\lambda}\gg1.
\end{eqn}
\end{frame}

\sectionslide{Ray series ansatz}

\begin{frame}
\centering
\includegraphics[height=0.55\paperheight]{L5F2_t00.pdf}

\raggedright
Start with a plane wave within the homogeneous left half-space:
\begin{eqn}
u_{i}(\mathbf{x},t)=f(t-p_{j}^{0}x_{j})a_{i}^{0}.
\end{eqn}
Can we extend this solution into the heterogeneous right half-space?
\end{frame}

\begin{frame}
\centering
\includegraphics[height=0.55\paperheight]{L5F2_t09.pdf}

\raggedright
As an initial guess, we look for a solution of the form:
\begin{eqn}
u_{i}(\mathbf{x},t)=f[t-T(\mathbf{x})]a_{i}(\mathbf{x}).
\end{eqn}
\end{frame}

\begin{frame}
Assumed solution takes the form:
\begin{eqn}
u_{i}(\mathbf{x},t)=f[t-T(\mathbf{x})]a_{i}(\mathbf{x}).
\end{eqn}
Introduction of the \textbf{travel time} function allows for curved wavefronts:
\begin{eqn}
T(\mathbf{x})=\text{constant}.
\end{eqn}
The \textbf{local slowness vector} is orthogonal to the wavefronts:
\begin{eqn}
p_{i}=\frac{\partial T}{\partial x_{i}}.
\end{eqn}
Polarisation vector allowed to vary with position -- accounts for \textbf{focusing effects} and for variations in
anisotropic properties with position.
\end{frame}

\begin{frame}
Harmonic decomposition solution:
\begin{eqn}
u_{i}(\mathbf{x},t)=\frac{1}{2\pi}\int_{\mathbb{R}}\tilde{f}(\omega)\ee^{\ii\omega[t-T(\mathbf{x})]}a_{i}(\mathbf{x})\dd\omega.
\end{eqn}
This is a convolution of $f$ with the singular solution:
\begin{eqn}
u_{i}(\mathbf{x},t)=\delta[t-T(\mathbf{x})]a_{i}(\mathbf{x}).
\end{eqn}
Experience has shown that we should generalise the singular part to:
\begin{eqn}
u_{i}(\mathbf{x},t)=\frac{1}{2\pi}\int_{\mathbb{R}}\sum_{n=0}^{\infty}\frac{a_{i}^{n}(\mathbf{x})}{(\ii\omega)^{n}}\ee^{\ii\omega[t-T(\mathbf{x})]}\dd\omega.
\end{eqn}
First term ($n=0$) in the \textbf{elastodynamic ray series} is the singular one, while successive ones become smoother
and smoother.
\end{frame}

\sectionslide{The eikonal equation}

\begin{frame}
Linearised equations with no force term:
\begin{eqn}
\rho\frac{\partial^{2}u_{i}}{\partial t^{2}}-\frac{\partial}{\partial x_{j}}\left(A_{ijkl}\frac{\partial u_{k}}{\partial x_{l}}\right)=0.
\end{eqn}
Transforming to the frequency domain:
\begin{eqn}
-\omega^{2}\rho u_{i}-\frac{\partial}{\partial x_{j}}\left(A_{ijkl}\frac{\partial u_{k}}{\partial x_{l}}\right)=0.
\end{eqn}
Ray series in the frequency domain:
\begin{eqn}
u_{i}(\mathbf{x},\omega)=\sum_{n=0}^{\infty}\frac{a_{i}^{n}(\mathbf{x})}{(\ii\omega)^{n}}\ee^{-\ii\omega T(\mathbf{x})}.
\end{eqn}
We can now substitute in our ansatz and see what happens!
\end{frame}

\begin{frame}
Example step from the calculation:
\begin{eqn}
\frac{\partial u_{k}}{\partial x_{l}}=\sum_{n=0}^{\infty}\frac{1}{(\ii\omega)^{n}}\left(\frac{\partial a_{k}^{n}}{\partial x_{l}}
-\ii\omega p_{l}a_{k}^{n}\right)\ee^{-\ii\omega T}.
\end{eqn}
A factor of frequency is added each time we differentiate the phase term.
\end{frame}

\begin{frame}
Equations of motion imply:
\begin{eqn}
\begin{aligned}
&\sum_{n=0}^{\infty}\frac{1}{(\ii\omega)^{n}}\Bigg[\frac{\partial}{\partial x_{j}}\left(A_{ijkl}\frac{\partial a_{k}^{n}}{\partial x_{l}}\right)
-\ii\omega\frac{\partial}{\partial x_{j}}(A_{ijkl}p_{l}a_{k}^{n})\\[\eqgap]
&\qquad-\ii\omega A_{ijkl}p_{j}\frac{\partial a_{k}^{n}}{\partial x_{l}}+(\ii\omega)^{2}A_{ijkl}p_{j}p_{l}a_{k}^{n}
-(\ii\omega)^{2}\rho a_{i}^{n}\Bigg]\ee^{-\ii\omega T}=0.
\end{aligned}
\end{eqn}
Common phase term cancels.

These equations must hold for \textbf{all frequencies} -- this implies that the \textbf{coefficients of each power of
frequency must separately vanish}.

Compare with power series solution of ODEs.
\end{frame}

\begin{frame}
\eqscale{20}
From the highest-order coefficient, we have:
\begin{eqn}
[A_{ijkl}p_{j}p_{l}-\rho\delta_{ik}]a_{k}^{0}=0.
\end{eqn}
At the next order:
\begin{eqn}
[A_{ijkl}p_{j}p_{l}-\rho\delta_{ik}]a_{k}^{1}=\frac{\partial}{\partial x_{j}}(A_{ijkl}p_{l}a_{k}^{0})+A_{ijkl}p_{j}\frac{\partial a_{k}^{0}}{\partial x_{l}}.
\end{eqn}
For the general term with $n\ge2$:
\begin{eqn}
\begin{aligned}
{}[A_{ijkl}p_{j}p_{l}-\rho\delta_{ik}]a_{k}^{n}={}&\frac{\partial}{\partial x_{j}}(A_{ijkl}p_{l}a_{k}^{n-1})+A_{ijkl}p_{j}\frac{\partial a_{k}^{n-1}}{\partial x_{l}}\\
&-\frac{\partial}{\partial x_{j}}\left(A_{ijkl}\frac{\partial a_{k}^{n-2}}{\partial x_{l}}\right).
\end{aligned}
\end{eqn}
First equation determines $\mathbf{p}$, while later ones successively determine the polarisation vectors at each order.

Where have we made an approximation?
\end{frame}

\begin{frame}
In this course we will consider only the equation derived from the \textbf{highest-order} term in the expansion.

If that is all that we require, we can systematically drop all but the highest-order terms within the previous derivation:
\begin{eqn}
\frac{\partial u_{k}}{\partial x_{l}}=\sum_{n=0}^{\infty}\frac{1}{(\ii\omega)^{n}}\left(\frac{\partial a_{k}^{n}}{\partial x_{l}}
-\ii\omega p_{l}a_{k}^{n}\right)\ee^{-\ii\omega T}\approx-\ii\omega p_{l}a_{k}^{0}\ee^{-\ii\omega T}.
\end{eqn}
This saves a lot of time!

You can practise these steps in the first problem set.
\end{frame}

\begin{frame}
From the highest-order coefficient, we have:
\begin{eqn}
[A_{ijkl}p_{j}p_{l}-\rho\delta_{ik}]a_{k}^{0}=0.
\end{eqn}
This motivates the definition of a \textbf{local Christoffel matrix}:
\begin{eqn}
\Gamma_{ik}(\mathbf{x},\mathbf{p})=\frac{1}{\rho(\mathbf{x})}A_{ijkl}(\mathbf{x})p_{j}p_{l}.
\end{eqn}
And so we can write the above condition as:
\begin{eqn}
\{\bm{\Gamma}[\mathbf{x},\mathbf{p}(\mathbf{x})]-\mathbf{I}\}\mathbf{a}^{0}(\mathbf{x})=\mathbf{0}.
\end{eqn}
Locally, things look identical to plane waves!
\end{frame}

\begin{frame}
The equation
\begin{eqn}
\{\bm{\Gamma}[\mathbf{x},\mathbf{p}(\mathbf{x})]-\mathbf{I}\}\mathbf{a}^{0}(\mathbf{x})=\mathbf{0}
\end{eqn}
has non-trivial solutions for the zeroth-order polarisation vector if and only if:
\begin{eqn}
\det\{\bm{\Gamma}[\mathbf{x},\mathbf{p}(\mathbf{x})]-\mathbf{I}\}=0.
\end{eqn}
This is a non-linear first-order PDE for the travel time called the \textbf{eikonal equation}.

If the local slowness surfaces are everywhere non-degenerate:
\begin{eqn}
\|\mathbf{p}(\mathbf{x})\|^{2}c_{k}[\mathbf{x},\hat{\mathbf{p}}(\mathbf{x})]^{2}-1=0.
\end{eqn}
\end{frame}

\sectionslide{The method of characteristics}

\begin{frame}
Recall the eikonal equation is:
\begin{eqn}
\|\mathbf{p}(\mathbf{x})\|^{2}c_{k}[\mathbf{x},\hat{\mathbf{p}}(\mathbf{x})]^{2}-1=0.
\end{eqn}
If we define the \textbf{ray Hamiltonian} for the $k$th sheet of the slowness surface:
\begin{eqn}
H_{k}(\mathbf{x},\mathbf{p})=\frac{1}{2}\|\mathbf{p}\|^{2}c_{k}(\mathbf{x},\hat{\mathbf{p}})^{2},
\end{eqn}
we can write the eikonal equation in the form:
\begin{eqn}
H_{k}[\mathbf{x},\mathbf{p}(\mathbf{x})]=\frac{1}{2}, \quad \mathbf{p}(\mathbf{x})=(\nabla T)(\mathbf{x}).
\end{eqn}
\end{frame}

\begin{frame}
\begin{columns}[c]
\column{0.5\textwidth}
\centering
\begin{tikzpicture}[line width=0.9pt,>={Stealth[length=3mm]},x=1.05cm,y=1.05cm]
\draw[->] (-0.3,0) -- (9.6,0) node[right] {$\mathbf{x}$};
\draw[->] (0,-0.3) -- (0,6.3) node[above] {$\mathbf{p}$};
\coordinate (P) at (2.8,3.0);
\draw[densely dotted] (2.8,0) -- (P);
\draw[dashed] (0,3.0) -- (P);
\node[below] at (2.8,-0.05) {$\mathbf{x}^{0}$};
\node[left] at (-0.05,3.0) {$\mathbf{p}^{0}$};
\draw (P) .. controls (4.6,3.1) and (5.4,4.3) .. (6.8,4.15) .. controls (7.8,4.05) and (8.4,4.25) .. (9.3,4.6);
\fill (P) circle (0.09);
\node[anchor=south west] (init) at (3.6,5.2) {initial data};
\draw[->,line width=0.6pt] (init.south west) ++(0.2,0) -- ([shift={(0.12,0.12)}]P);
\node[align=left,anchor=north west] (lev) at (5.4,2.6) {solution, within\\$H_{k}(\mathbf{x},\mathbf{p})=\tfrac{1}{2}$};
\draw[->,line width=0.6pt] (lev.north) ++(0.3,0) -- (6.3,4.05);
\end{tikzpicture}
\column{0.46\textwidth}
Solution of the eikonal equation is a 3D surface lying within a level set of the ray Hamiltonian in a 6D
phase space.

Initial data specifies a 2D surface within this level set.

The problem is how to extend this 2D surface into the appropriate 3D surface.

This is where Hamilton's equations enter into the method.
\end{columns}
\end{frame}

\begin{frame}
Initial data extended through integration of Hamilton's equations, with $\sigma$ a \textbf{generating parameter}:
\begin{eqn}
\begin{gathered}
\frac{\ddns x_{i}}{\ddns\sigma}=\frac{\partial H_{k}}{\partial p_{i}}, \quad \frac{\ddns p_{i}}{\ddns\sigma}=-\frac{\partial H_{k}}{\partial x_{i}},\\[\eqgap]
x_{i}(0)=x_{i}^{0}, \quad p_{i}(0)=p_{i}^{0}.
\end{gathered}
\end{eqn}
This works because the Hamiltonian is constant along the integral curves:
\begin{eqn}
\begin{aligned}
\frac{\ddns}{\ddns\sigma}H_{k}[\mathbf{x}(\sigma), \mathbf{p}(\sigma)]&=\frac{\partial H_{k}}{\partial x_{i}}\frac{\ddns x_{i}}{\ddns\sigma}+\frac{\partial H_{k}}{\partial p_{i}}\frac{\ddns p_{i}}{\ddns\sigma}\\[\eqgap]
&=-\frac{\ddns p_{i}}{\ddns\sigma}\frac{\ddns x_{i}}{\ddns\sigma}+\frac{\ddns x_{i}}{\ddns\sigma}\frac{\ddns p_{i}}{\ddns\sigma}=0.
\end{aligned}
\end{eqn}
There are a few subtler points to the method that are being ignored.
\end{frame}

\begin{frame}
\textbf{Rays} are the projection of integral curves into the position domain.

Rate of change of the travel time along a ray:
\begin{eqn}
\frac{\ddns}{\ddns\sigma}T[\mathbf{x}(\sigma)]=\frac{\partial T}{\partial x_{i}}\frac{\ddns x_{i}}{\ddns \sigma}=p_{i}\frac{\partial H_{k}}{\partial p_{i}}.
\end{eqn}
Right-hand side simplified using Euler's theorem for homogeneous functions (proof in the handout, non-examinable):
\begin{eqn}
p_{i}\frac{\partial H_{k}}{\partial p_{i}}=2H_{k}\arrowto\frac{\ddns }{\ddns \sigma}T[\mathbf{x}(\sigma)]=1.
\end{eqn}
\textbf{The generating parameter along the ray is equal to the travel time.}
\end{frame}

\begin{frame}
\centering
\includegraphics[height=0.6\paperheight]{L5F3.pdf}

\raggedright
If rays for two different integral curves cross, then the travel time need not be well defined. The envelope of the
ray family is known as a \textbf{caustic}. The \textbf{ray theoretic amplitude is infinite} there, while the true
amplitude is large but finite.
\end{frame}

\begin{frame}
\centering
\includegraphics[width=0.9\textwidth]{L5F4.pdf}

\raggedright
Within heterogeneous materials, caustics occur frequently -- cf.\ Penrose--Hawking singularity theorems in general
relativity.
\end{frame}

\begin{frame}[t]
\heading{What you need to know and be able to do}
\begin{enumerate}
\item[(i)] The form of the ray-series expansion, and understand its underlying physical motivation.
\item[(ii)] How to obtain the eikonal equation and transport equations by substituting the ray expansion into the
equations of motion. The full derivation given in the lecture is too involved to be set in an exam, but you may be
asked to apply these methods within simpler situations.
\item[(iii)] Understand how the eikonal equation can be interpreted geometrically within a six-dimensional phase
space, and be able to explain how solutions can be obtained through integration of Hamilton's equations.
\item[(iv)] Be able to solve the Hamiltonian ray equations in simple cases, and also to deduce some of their
properties. Practice in this will be given in the next lecture and within the first problem set.
\item[(v)] What a caustic is, and why ray theory breaks down there.
\end{enumerate}
\end{frame}

\end{document}
